\chapter{Bevezetés}

Diplomamunkám a \textit{SuperSet} névre keresztelt, általam fejlesztett, elosztott fordítóprogram teszt keretrendszert
mutatja be. A fordítóprogramok tesztelése egy korántsem triviális téma. Különböző módszertanok a fordítóprogramok
különböző aspektusait, kimeneteit hivatottak tesztelni. A \textit{SuperSet} célja egy olyan teszt keretrendszert
nyújtani, ami az alábbi tulajdonságokkal rendelkezik:

\begin{itemize}
    \item Támogatja a nyelvi \emph{konformanciát} igazoló, valamint \emph{fuzz} tesztek futtatását.
    \item Könnyen, de mégis részletesen és robusztusan konfigurálható. Lehetővé teszi fordítóprogram konfigurációinak
          kombinációját és kizárását is.
    \item Lehetővé teszi különböző tesztesetek, tesztecsomagba való szervezését.
    \item Keresztfordítók esetén képes a futásidejű teszteket szimulátor alkalmazásokban futtatni.
    \item A tesztcsomag szerzőnek lehetőséget biztosít a konkrét fordítóprogram invokációt \\ programozottan
          változtatni egy bővítmény rendszeren keresztül.
    \item A kimenete \textit{JUnit XML} formátumú, hogy egyéb vizualizációs eszközök fel tudják dolgozni azt.
    \item Nem utolsó sorban lehetőséget ad az elosztott futtatásra, hogy a nagy méretű tesztcsomagok gyorsabban
          fussanak feltéve, hogy extra hardver áll rendelkezésre.
\end{itemize}

A projekt érdekessége, hogy a Python programozási nyelv egy relatíve új, extra típus annotációkkal ellátott
változatában készült. A legtöbb típus annotációnak nincs futásidejű szemantikája, ugyanakkor ellenőrizhetőek egy
harmadik féltől származó programokkal~\cites{mypy,pyright,ty}. A típusosság nagyban megkönnyítette a fejlesztési
folyamatot és nagyobb konfidenciát adott a kód helyességére nézve. Az implementáció során törekedtem a
kódolvashatóságra, a \textit{DRY}\footnote{\textit{\foreignlanguage{american}{Don't repeat yourself}}, a kód
tömörségére érdemes törökedni.} és a \textit{SOLID}\footnote{\textit{SOLID}, az objektum orientált programozásban öt
tervezési szabály.} programozási elvek betartására.%
